\documentclass[11pt]{article}
\usepackage{manual_docs_style}
\externaldocument{build/terminology_shared_utilities}[terminology_shared_utilities.pdf]
\externaldocument{build/appendix_environment}[appendix_environment.pdf]

\begin{document}

\docTitle{Environments}
\phantomsection\label{docstart:environment_profiles}

\section*{Overview}

This document describes the manually maintained files under \code{models/simulink/<simulator_id>/} that define a \name{} \termSimulator.
A \termSimulator becomes usable once the required files exist, \code{build_metadata.py} has been run for \code{simulink_model_original.mdl}, and the \termSimulator has been added to \code{shared/benchmark_utils.py::ALLOWED_TSENV_MODELS}.

\phantomsection\label{sec:allowed-tsenv-models}
The files are located in \code{models/simulink/<simulator_id>/}.
\begin{DocCode}
models/simulink/<simulator_id>/
|-- simulink_model_original.mdl
|-- build_metadata.py
|-- description_levels.json
|-- experiment_config.json
|-- noise_adder.py
|-- generation_policies/<policy_id>.json
|-- detectability_specific_environment.py  optional
|-- features.py                            optional
`-- basic_rule.py                          optional
\end{DocCode}

\section*{Generation Policy Schema}
\phantomsection\label{sec:generation-policy-schema}

A generation policy defines:

\begin{itemize}
\item how many candidate families to generate;
\item how to sample baseline parameters and initial states;
\item how to sample intervention times;
\item how to sample post-intervention values;
\item how many children to generate per root-cause parameter and direction.
\end{itemize}

Candidate generation resolves the policy into a concrete plan before simulation.

\subsection*{Top-Level Fields}

Required top-level fields in a generation policy. All fields listed in this table are required.

\begin{center}
\begin{tabular}{p{0.28\linewidth}p{0.60\linewidth}}
\hline
\textbf{Field} & \textbf{Meaning} \\
\hline
\code{policy_id} & Non-empty identifier for the policy. Should match the filename stem. \\
\code{model} & \termSimulator identifier under \code{models/simulink/<simulator_id>}. \\
\code{seed} & Integer seed used to make candidate generation reproducible. \\
\code{target_candidate_families} & Number of baseline-centered candidate families to sample. \\
\code{baseline_sampler} & Defines how baseline parameters and initial states are sampled. \\
\code{intervention_sampler} & Defines intervention parameters, directions, times, and values. \\
\hline
\end{tabular}
\end{center}

\subsection*{Minimal Shape}

\begin{DocCode}
{
  "policy_id": "production_v1",
  "model": "BallDrop",
  "seed": 456,
  "target_candidate_families": 1000000,
  "baseline_sampler": {
    "method": "latin_hypercube"
  },
  "intervention_sampler": {
    "parameters": "all",
    "directions": ["increase", "decrease"],
    "children_per_parameter_direction": 3,
    "intervention_time": {
      "distribution": "uniform",
      "min_fraction_of_horizon": 0.2,
      "max_fraction_of_horizon": 0.8
    },
    "value_sampler": {
      "method": "conditional_uniform",
      "min_abs_distance_multiplier": 1.0,
      "min_srd_distance_multiplier": 1.0
    }
  }
}
\end{DocCode}

\subsection*{\texorpdfstring{\titlecode{baseline_sampler}}{baseline\_sampler}}

\code{baseline_sampler} defines how baseline parameters and initial-state variables are sampled.
It samples over the variables listed in:

\begin{DocCode}
experiment_config.json::exposed_variables.parameters
experiment_config.json::exposed_variables.initial_state
\end{DocCode}

The sampler must respect all \code{allowed_intervals} defined in \code{experiment_config.json}.

\subsubsection*{Schema}

\begin{DocCode}
"baseline_sampler": {
  "method": "latin_hypercube",
  "variables": {
    "mass": {
      "distribution": "loguniform"
    },
    "gravity": {
      "distribution": "uniform"
    },
    "initial_height": {
      "distribution": "uniform"
    }
  },
  "baseline_min_distance_scale": 0.2
}
\end{DocCode}

\subsubsection*{Fields}

\begin{itemize}
\item \code{method}: required string. Allowed values are \code{independent}, \code{latin_hypercube}, \code{sobol}, and \code{grid}.
\item \code{variables}: optional object keyed by exposed variable name. It overrides the default distribution for selected variables.
\item \code{baseline_min_distance_scale}: optional non-negative number. If present, generated baselines should be separated from one another according to the policy's diversity constraint.
\end{itemize}

Allowed variable-level distributions are:

\begin{itemize}
\item \code{uniform}: sample uniformly in the legal interval.
\item \code{loguniform}: sample log-uniformly in the legal interval. The legal interval must be strictly positive.
\item \code{boundary_biased}: sample with extra mass near the lower and upper legal bounds.
\item \code{fixed}: use a fixed value. Requires a \code{value} field.
\end{itemize}

A fixed value must still lie inside the variable's \code{allowed_intervals}.

\subsection*{\texorpdfstring{\titlecode{intervention_sampler}}{intervention\_sampler}}

\code{intervention_sampler} defines how intervention children are generated from each baseline family.

\subsubsection*{Schema}

\begin{DocCode}
"intervention_sampler": {
  "parameters": "all",
  "directions": ["increase", "decrease"],
  "children_per_parameter_direction": 3,
  "intervention_time": {
    "distribution": "uniform",
    "min_fraction_of_horizon": 0.2,
    "max_fraction_of_horizon": 0.8
  },
  "value_sampler": {
    "method": "conditional_uniform",
    "min_abs_distance_multiplier": 1.0,
    "min_srd_distance_multiplier": 1.0
  }
}
\end{DocCode}

\subsubsection*{Fields}

\begin{itemize}
\item \code{parameters}: either \code{"all"} or a non-empty list of parameter names from \code{experiment_config.json::exposed_variables.parameters}.
\item \code{directions}: non-empty list containing \code{"increase"}, \code{"decrease"}, or both.
\item \code{children_per_parameter_direction}: positive integer.
\item \code{intervention_time}: object defining how intervention times are sampled.
\item \code{value_sampler}: object defining how post-intervention values are sampled.
\end{itemize}

Only variables listed in \code{experiment_config.json::exposed_variables.parameters} may be intervened on.
Variables listed under \code{initial_state} are baseline-only variables and must not change after simulation start.

\subsubsection*{\texorpdfstring{\titlecode{intervention_time}}{intervention\_time}}

\code{intervention_time} defines how the step-change time is sampled.

\begin{DocCode}
"intervention_time": {
  "distribution": "uniform",
  "min_fraction_of_horizon": 0.2,
  "max_fraction_of_horizon": 0.8
}
\end{DocCode}

Allowed fields:

\begin{itemize}
\item \code{distribution}: required string. Allowed values are \code{uniform}, \code{fixed}, and \code{choice}.
\item \code{min_fraction_of_horizon}: required for \code{uniform}. Float in \code{[0,1]}.
\item \code{max_fraction_of_horizon}: required for \code{uniform}. Float in \code{[0,1]}.
\item \code{value}: required when \code{distribution} is \code{fixed}.
\item \code{values}: required when \code{distribution} is \code{choice}.
\end{itemize}

All intervention times must satisfy:

\begin{DocCode}
0 < intervention_time < experiment_config.json::end_time_input_s
\end{DocCode}

If \code{min_fraction_of_horizon} and \code{max_fraction_of_horizon} are used, the actual interval is computed from \code{experiment_config.json::end_time_input_s}.

\subsubsection*{\texorpdfstring{\titlecode{value_sampler}}{value\_sampler}}

\code{value_sampler} defines how the post-intervention value is sampled for each changed parameter.

\begin{DocCode}
"value_sampler": {
  "method": "conditional_uniform",
  "min_abs_distance_multiplier": 1.0,
  "min_srd_distance_multiplier": 1.0
}
\end{DocCode}

Allowed \code{method} values:

\begin{itemize}
\item \code{conditional_uniform}: uniformly sample valid values on the requested side of the baseline value after applying distance constraints.
\item \code{boundary_biased}: sample valid values with extra mass near the legal interval boundaries.
\item \code{fixed_delta}: add or subtract a fixed amount. Requires \code{delta}.
\item \code{fixed_srd}: sample values targeting a fixed SRD distance. Requires \code{target_srd}.
\end{itemize}

The generated post-intervention value must satisfy all of the following:

\begin{itemize}
\item it lies inside \code{experiment_config.json::exposed_variables.parameters::<parameter>::allowed_intervals};
\item it changes exactly one root-cause parameter;
\item its absolute distance from the baseline value is at least the environment-level \termMinAbsDist{} multiplied by \code{min_abs_distance_multiplier};
\item its SRD distance from the baseline value is at least the environment-level \termMinSRDDistance{} multiplied by \code{min_srd_distance_multiplier};
\item its direction agrees with the requested direction.
\end{itemize}

If no valid value exists for a requested parameter and direction, the generator must not create a runnable child node.
The missing direction must instead be recorded in \code{generation_report.json}.

\section*{\texorpdfstring{\titlecode{description_levels.json}}{description\_levels.json}}

\code{description_levels.json} contains \termSimulator-specific text used by the prompt renderer.
It has:
\begin{itemize}
\item \code{textual_description}: object keyed by \termDescLevel.
\begin{itemize}
\item \code{high}: detailed natural-language description, including relevant equations, observed channels, and root-cause parameters.
\end{itemize}
\item \code{observed_signals_description}: signal description used in prompts.
\item \code{internal_naming_to_agent_facing_signal}: object mapping internal signal names to agent-facing signal descriptions.
The non-\code{time} keys must match \code{experiment_config.json::observable_signals::observable_signals} exactly and in the same order.
\begin{DocCode}
{
  "signal_internal_name_0": "agent-facing description 0",
  "signal_internal_name_1": "agent-facing description 1",
  "time": "time"
}
\end{DocCode}
\item \code{internal_naming_to_agent_facing_parameter}: object mapping internal intervenable parameter names from \code{experiment_config.json::exposed_variables.parameters} to agent-facing labels.
\end{itemize}


\section*{\texorpdfstring{\titlecode{experiment_config.json}}{experiment\_config.json}}

\code{experiment_config.json} defines the environment contract.
It should contain:

\begin{itemize}
\item \code{paper_facing_name}: environment name used in the paper and generated tables.
\item \code{exposed_variables}: variables externally exposed to \name{}.
\item \code{observable_signals}: non-time channels observed by agents and metrics.
\item \code{sampling_rate_hz}: target sampling rate for materialized trajectories.
\item \code{end_time_input_s}: simulation horizon.
\item \code{detectability}: thresholds used by Stage: Cheap Filter Metrics.
\end{itemize}

\subsection*{\code{exposed_variables}}

\code{exposed_variables} has two groups:

\begin{itemize}
\item \code{parameters}: root-cause parameters that may be changed by an intervention.
\item \code{initial_state}: baseline-only state variables that are sampled at the start of a run and must not change after simulation start.
\end{itemize}

Each exposed variable should define:

\begin{itemize}
\item \code{allowed_intervals}: legal range \code{[min, max]}.
\item \termMinSRDDistance
\item \termMinAbsDist
\end{itemize}

\subsection*{\code{observable_signals}}

\code{observable_signals} defines the channels that agents and downstream metrics see.

\begin{itemize}
\item \code{observable_signals}: ordered list of non-time signal names.
\item \code{signal_type}: object keyed by signal name. Each value has \code{type}, which is either \code{continuous} or \code{impulse_like}.
\end{itemize}

\subsection*{\code{detectability}}

\code{detectability} contains the thresholds used by Stage: Cheap Filter Metrics to decide whether an intervention is visible in the observed signals.
It is divided into continuous-signal rules, impulse-like-signal rules, and per-signal RMS thresholds.

\begingroup
\footnotesize
\renewcommand{\arraystretch}{1.12}
\setlength{\tabcolsep}{4pt}
\begin{longtable}{|P{0.20\textwidth}|P{0.28\textwidth}|P{0.42\textwidth}|}
\caption{Detectability fields.}\\
\hline
\textbf{Block} & \textbf{Field} & \textbf{Meaning} \\
\hline
\endfirsthead
\hline
\textbf{Block} & \textbf{Field} & \textbf{Meaning} \\
\hline
\endhead
\code{continuous} & \code{min_srd_distance} & Minimum \termSRD distance used for threshold-crossing detection on continuous signals. \\
\hline
\code{continuous} & \code{epsilon_SRD} & Epsilon term used in the \termSRD denominator. \\
\hline
\code{continuous} & \code{minimum_consecutive_below_SRD} & Number of consecutive samples that must exceed the \termSRD threshold before a difference is recorded. \\
\hline
\code{impulse_like} & \code{min_srd_distance} & Minimum \termSRD distance used for impulse-like signals. \\
\hline
\code{impulse_like} & \code{epsilon_SRD} & Epsilon term used in the \termSRD denominator for impulse-like signals. \\
\hline
\code{RMS_thresholds} & \code{<signal_name>} & Per-signal RMS threshold used to decide whether two trajectories are sufficiently different. \\
\hline
\end{longtable}
\endgroup

\section*{\texorpdfstring{\titlecode{noise_adder.py}}{noise\_adder.py}}

\code{noise_adder.py} implements the \termSimulator-specific observation-noise profiles.
It must expose:

\begin{itemize}
\item \code{NOISE_DICT}: \termSimulator-specific base observation-noise magnitudes for the low-noise regime, typically keyed by observable signal. The high-noise regime is obtained by multiplying these magnitudes by the shared \code{NOISE_LEVEL_MULTIPLIER["high"]}; \code{none} uses multiplier \code{0}.
\phantomsection\label{def:noise-adder-interface}
\end{itemize}

\begin{DocCode}
def quantify_noise(clean: pd.DataFrame, noisy: pd.DataFrame, reference: pd.DataFrame) -> dict
\end{DocCode}

\begin{DocCode}
def add_noise(src: pd.DataFrame, seed: int, noise_level: Literal["low", "high"], ref: pd.DataFrame | None = None) -> tuple[pd.DataFrame, dict]
\end{DocCode}

\section*{Optional Implementation Hooks}

\code{detectability_specific_environment.py} may expose:

\begin{DocCode}
def is_detectable(intervention: pd.DataFrame, baseline: pd.DataFrame, run_parameters: dict, intervention_time: float, min_first_diff: float | None) -> str
\end{DocCode}

This hook is an environment-specific veto used by Stage: Cheap Filter Metrics.
If it returns \code{"no"}, the child is marked as not detectable even if the generic RMS/SRD checks pass.
\subsection*{Optional files}

\begin{itemize}
\item \code{features.py}: A self-contained script that contains a feature extractions set. It has a \code{FEATURE_NAMES: tuple[str, ...]} with a list of all available features.
Each feature can be exposed independently. It also has utility functions that are shared among feature. 
It's used by basic \code{rule.py} and by \code{detectability_specific_environment.py}.
\item \code{basic_rule.py}: Optional baseline prediction rule used by metric-computation sweeps.
\end{itemize}

\phantomsection\label{def:basic-rule-predict-interface}
If present, \code{basic_rule.py} must expose:

\begin{DocCode}
def predict(df) -> list[str]
\end{DocCode}


\section*{Validation Rules}
\phantomsection\label{sec:environment-validation-rules}

These checks happen before simulation or candidate generation:

\begin{itemize}
\item \code{description_levels.json::textual_description::<desc_level>} for \termDescLevel in \{\code{high}\} must start with ``by simulating'' when it is appended to the shared prompt text.
\item \code{description_levels.json::internal_naming_to_agent_facing_signal} must be a non-empty object.
\item The keys of \code{description_levels.json::internal_naming_to_agent_facing_signal}, excluding \code{time}, must match \code{experiment_config.json::observable_signals::observable_signals} exactly and in the same order.
\item The signals listed in \code{experiment_config.json::observable_signals::observable_signals} must appear in either \code{generated/metadata.json::simulink_signals_available} or \code{generated/metadata.json::simscape_signals_available}.
\end{itemize}

\paragraph{Modeling requirement.}
Benchmark samples should be generated from physically valid simulator runs and filtered in two stages.
First, the cheap-filter stage rejects failed, numerically invalid, low-SNR, and non-detectable samples using the matched baseline, the matched static-change probe, and any environment-specific veto hooks.
Second, for benchmark selection, the surrogate confidence filter keeps only samples for which a calibrated neural surrogate assigns high confidence to the simulator ground-truth label.
This confidence gate is a scalable recognizability filter, not a proof of causal uniqueness.
The optional Oracle Audit may be run on a subset for post-hoc ambiguity analysis, but it is not part of default selection.

The full environment-by-environment catalog is collected in the paper appendix, under \hyperref[docstart:appendix_environment]{Appendix: Environment Catalog}.

\end{document}
